%% SoftwareX Original Software Publication manuscript.
%% Prepared from the official SoftwareX OSP LaTeX structure (accessed 2026-07-15).
\documentclass[preprint,12pt,a4paper]{elsarticle}

\usepackage{amssymb}
\usepackage{amsmath}
\usepackage{booktabs}
\usepackage{array}
\usepackage{hyperref}
\usepackage{url}

\hypersetup{
  colorlinks=true,
  linkcolor=blue!55!black,
  citecolor=blue!55!black,
  urlcolor=blue!55!black
}
\urlstyle{same}
\setlength{\parindent}{0pt}
\journal{SoftwareX}

\begin{document}

\begin{frontmatter}

\title{Agent Wiki: A local-first evidence-closed Markdown knowledge base for LLM agents}

\author[independent]{Zijun Chu\corref{cor1}}
\ead{chunimav5@gmail.com}
\ead[url]{https://github.com/NimaChu/agent-wiki}
\cortext[cor1]{Corresponding author.}
\address[independent]{Independent Researcher, Room 1303, No. 19, Lane 899,
Xinfeng Middle Road, Qingpu District, Shanghai, China}

\begin{abstract}
Agent Wiki is an open-source, local-first knowledge substrate for large language
model agents. It separates captured evidence from synthesized Markdown wiki
pages, resolves ordinary wiki links as a computable graph, and checks an
evidence-closure rule before a source is treated as processed. The software also
supports source capture, vault search, maintenance queues, link repair, image
evidence inventories, optional external-source import, and a read-only graph
dashboard. A technical-documentation case study containing 1,137 raw source
notes and 127 wiki pages produced 5,444 graph edges with no unresolved links,
orphaned wiki pages, or processed-source closure violations under the included
lint checks. Agent Wiki provides an inspectable corpus-preparation and
maintenance layer that can be used independently or upstream of retrieval-
augmented generation.
\end{abstract}

\begin{keyword}
LLM agents \sep knowledge management \sep Markdown \sep local-first software
\sep provenance \sep software engineering
\end{keyword}

\end{frontmatter}

\section*{Required Metadata}

\section*{Current code version}

\begin{table}[!h]
\small
\begin{tabular}{|p{0.7cm}|p{5.1cm}|p{7.2cm}|}
\hline
\textbf{Nr.} & \textbf{Code metadata description} & \textbf{Current Agent Wiki release} \\
\hline
C1 & Current code version & v0.2.1 \\
\hline
C2 & Permanent GitHub link to code/repository used for this code version &
\url{https://github.com/NimaChu/agent-wiki/releases/tag/v0.2.1} \\
\hline
C3 & Legal Code License & MIT License \\
\hline
C4 & Code versioning system used & git \\
\hline
C5 & Software code languages, tools, and services used &
JavaScript (ECMAScript modules), Node.js, React, Vite, PowerShell \\
\hline
C6 & Compilation requirements, operating environments \& dependencies &
Node.js 18 or newer and npm. The command-line workflows are file-based; the
optional dashboard requires a modern web browser. Validated on Windows 11 with
Node.js 20.20.0. \\
\hline
C7 & Link to developer documentation/manual &
\url{https://github.com/NimaChu/agent-wiki/blob/v0.2.1/README.md} \\
\hline
C8 & Support email for questions & \href{mailto:chunimav5@gmail.com}{chunimav5@gmail.com} \\
\hline
\end{tabular}
\caption{Code metadata (mandatory).}
\label{tab:code-metadata}
\end{table}

\section{Motivation and significance}

Large language model (LLM) agents now perform multi-step reading, coding, and
maintenance tasks. These workflows need durable knowledge: source material,
project conventions, decisions, screenshots, and unresolved follow-ups must
remain available after a chat session ends. Existing approaches usually rely
on chat history, human-authored note applications, application-specific memory,
or retrieval indexes. Each is useful, but independent developers and small
research teams still lack a simple way to maintain source-grounded agent
knowledge that is private, inspectable, and executable with ordinary developer
tools.

Agent Wiki addresses this problem with a local-first repository rather than a
hosted database. Local-first software emphasizes user ownership and continued
access to data independent of a cloud service [1]. Agent Wiki applies that
principle to agent-maintained knowledge: Markdown files are the source of truth,
while graph data, dashboards, and future retrieval indexes are derived
artifacts. Users can inspect, diff, search, version, or remove the knowledge
base without a proprietary export process.

The software is complementary to retrieval-augmented generation (RAG) [2]. RAG
optimizes answer-time retrieval, whereas Agent Wiki prepares and audits the
corpus before retrieval. It preserves captured evidence, synthesizes reusable
knowledge units, and makes maintenance state computable. This is particularly
relevant to tool-using and software-engineering agents, which interact with
repositories and terminals over extended workflows [3,4].

The intended users are researchers, developers, and technically curious
individuals who want an agent-maintained knowledge base without deploying a
database or embedding service. A user clones the repository, installs npm
dependencies, and asks a coding agent to operate the included capture,
maintenance, search, and validation commands. The same repository rules are
readable by both humans and agents.

\section{Software description}

\subsection{Software architecture}

Agent Wiki has two persistent knowledge layers and three derived service
layers. The raw layer stores source captures, provenance metadata, snapshots,
content hashes, image inventories, and processing state. The wiki layer stores
atomic, durable pages for concepts, entities, methods, APIs, workflows,
comparisons, and recurring questions. Wiki pages link back to raw evidence with
Obsidian-compatible links.

The command layer scans both persistent layers and implements capture, search,
status, lint, maintenance queues, link repair, image extraction, and optional
external-source import. The graph layer resolves links by path, title, basename,
and alias, then derives nodes, edges, unresolved targets, relation hints, and
incoming/outgoing counts. The optional React/Vite dashboard reads generated
JSON and never writes to the vault. Table~\ref{tab:architecture} summarizes the
architecture without introducing a runtime database.

\begin{table}[!h]
\small
\centering
\begin{tabular}{p{3.0cm}p{4.6cm}p{6.2cm}}
\toprule
Component & Primary representation & Responsibility \\
\midrule
Raw evidence & Markdown, snapshots, images & Preserve provenance and source material. \\
Wiki knowledge & Atomic Markdown pages & Synthesize reusable, source-linked knowledge. \\
Local CLI & Node.js scripts & Capture, search, maintain, lint, and repair the vault. \\
Derived graph & In-memory scan and JSON & Resolve links and compute health metrics. \\
Dashboard & React/Vite web application & Provide read-only graph and evidence inspection. \\
\bottomrule
\end{tabular}
\caption{Agent Wiki architecture and ownership boundaries.}
\label{tab:architecture}
\end{table}

\subsection{Software functionalities}

The central maintenance mechanism is evidence closure. Let $r$ be a raw source,
$T(r)$ its declared wiki targets, $R(t)$ the link resolver, and $L(w)$ the links
from a wiki page. A processed source must satisfy

\[
|T(r)|>0 \land \forall t\in T(r):R(t)\neq\bot \land
\exists w\in R(T(r)):r\in L(w) \land \neg\mathrm{followup}(r).
\]

The lint command reports missing targets, unresolved targets, missing wiki
backlinks, and explicit follow-up flags. It also reports unresolved links,
orphaned wiki pages, invalid relation hints, weak pages, and missing metadata.
The invariant does not prove semantic correctness; it makes the connection
between processed sources and synthesized knowledge auditable.

For visual sources, the image workflow extracts Markdown and HTML image
references, resolves relative URLs, mirrors remote images, writes an
\texttt{image-index.json}, records context and download failures, and inserts a
compact inventory into the raw note. Selected screenshots or diagrams can then
be promoted into wiki pages while the complete inventory remains in the source
layer.

\section{Illustrative examples}

A minimal installation and health check is:

\begin{verbatim}
git clone https://github.com/NimaChu/agent-wiki.git
cd agent-wiki
npm install
npm run wiki:status
npm run wiki:lint
\end{verbatim}

A user can capture a web source, search the resulting vault, and inspect the
optional graph with:

\begin{verbatim}
npm run wiki:capture -- --title "Source title" \
  --url "https://example.com"
npm run wiki:search -- "query terms"
npm run dashboard:open
\end{verbatim}

In normal use, a coding agent reads the repository rules, places captured
material in \texttt{raw/}, distills reusable concepts into \texttt{wiki/},
links claims to evidence, and runs lint before marking a source processed. The
dashboard is generated only when requested; it is not required for capture,
maintenance, or query workflows.

As an illustrative case study, Agent Wiki maintained a local technical-
documentation vault centered on Autodesk FlexSim material. On 2026-07-15,
\texttt{wiki:status} reported 1,275 nodes, 5,444 edges, 1,137 raw sources, and
127 wiki pages. Of the raw sources, 1,136 were processed and none were pending,
stale, or marked for follow-up. \texttt{wiki:lint} reported no unresolved links,
invalid relations, processed-source issues, orphaned or weak wiki pages,
missing claim sources, or required metadata omissions. This result measures
mechanical graph health, not factual answer quality.

\section{Impact}

Agent Wiki changes the unit of agent memory from an opaque application state to
a reviewable software artifact. This has three practical effects. First,
provenance survives synthesis: a concept page can be traced to one or more raw
notes and locally mirrored visual evidence. Second, maintenance becomes
testable: a user or continuous-integration job can run the same status and lint
commands as the agent. Third, the knowledge base remains reusable: Markdown can
feed text search, graph navigation, embeddings, RAG, or future Model Context
Protocol services without changing the source-of-truth format.

The case study provides initial evidence that the design can maintain link and
closure health across more than one thousand source notes. Its value is not the
node count alone but the absence of unresolved links, orphaned wiki pages, and
processed-source violations under executable checks. These properties support
research workflows in which literature, documentation, protocols, or software
knowledge must remain inspectable while agents update the synthesis layer.

The software is domain independent. Its repository conventions can be applied
to technical documentation, literature reviews, laboratory procedures, policy
collections, or project knowledge. The raw/wiki separation also supports a
privacy-preserving publication model: reusable tooling can be shared publicly
while private or bulky knowledge remains outside version control.

Current adoption evidence is limited to a single-developer, single-corpus case
study. Agent Wiki has not yet been evaluated through a controlled user study or
against chat-only and RAG-first baselines. Future work should measure answer
quality, maintenance effort, cross-agent consistency, and user trust; extend the
public case-study snapshot with semantic annotations; and evaluate optional
claim-level citations and embedding indexes. The public software distribution and archived release [5]
allow these questions to be pursued without depending on a hosted service.

\section{Conclusions}

Agent Wiki provides a local-first, Markdown-native substrate for durable LLM-
agent knowledge. Its principal software contribution is an executable workflow
that separates raw evidence from synthesized wiki pages and checks evidence
closure before sources are treated as processed. Capture, search, lint, image
evidence, graph generation, and dashboard inspection are available through a
small set of local commands. The documentation case study demonstrates
mechanical health at more than one thousand raw sources while also defining the
limits of the present evidence. Agent Wiki is intended as an inspectable layer
upstream of, or independent from, retrieval systems.

\section*{Current executable software version}

\begin{table}[!h]
\small
\begin{tabular}{|p{0.7cm}|p{5.1cm}|p{7.2cm}|}
\hline
\textbf{Nr.} & \textbf{Executable software metadata description} & \textbf{Current Agent Wiki release} \\
\hline
S1 & Current software version & v0.2.1 \\
\hline
S2 & Permanent link to executables of this version &
\url{https://github.com/NimaChu/agent-wiki/releases/tag/v0.2.1} \\
\hline
S3 & Permanent link to reproducible capsule/archive &
\url{https://doi.org/10.5281/zenodo.21252445} \\
\hline
S4 & Legal Software License & MIT License \\
\hline
S5 & Computing platforms/operating systems &
Node.js 18+ environment; validated on Windows 11. \\
\hline
S6 & Installation requirements and dependencies & Node.js 18+ and npm; run \texttt{npm install}. \\
\hline
S7 & Link to user manual &
\url{https://github.com/NimaChu/agent-wiki/blob/v0.2.1/README.md} \\
\hline
S8 & Support email for questions & \href{mailto:chunimav5@gmail.com}{chunimav5@gmail.com} \\
\hline
\end{tabular}
\caption{Executable software metadata (optional).}
\label{tab:software-metadata}
\end{table}

\section*{Data availability}

The reusable software and manuscript materials are available from GitHub and
Zenodo [5]. A public, derived FlexSim 2026 case-study subset is included in the
release at \url{https://github.com/NimaChu/agent-wiki/tree/v0.2.1/datasets/flexsim-2026-case-study}.
It contains metadata and graph structure for 1,092 source records and 88
synthesized wiki records (1,180 nodes and 5,163 resolved edges), but excludes
captured Autodesk text, snapshots, and images. The manuscript's full-vault
counts also include non-FlexSim local sources and pages that are not public.

\section*{CRediT author statement}

Zijun Chu: Conceptualization, Methodology, Software, Validation, Investigation,
Data curation, Writing--original draft, Writing--review and editing,
Visualization, Project administration.

\section*{Declaration of competing interest}

The author declares no known competing financial interests or personal
relationships that could have appeared to influence the work reported in this
paper.

\section*{Funding}

This research did not receive any specific grant from funding agencies in the
public, commercial, or not-for-profit sectors.

\section*{Declaration of generative AI and AI-assisted technologies in the manuscript preparation process}

During the preparation of this work, the author used OpenAI Codex to assist with
manuscript organization, language drafting, and LaTeX preparation. The author
reviewed and edited the content as needed and takes full responsibility for the
content of the article.

\section*{Acknowledgements}

The author thanks the open-source communities whose tools and documentation
support local-first software development.

\begin{thebibliography}{00}

\bibitem{kleppmann2019localfirst}
Kleppmann M, Wiggins A, van Hardenberg P, McGranaghan M. Local-first software:
You own your data, in spite of the cloud. In: Proc ACM SIGPLAN Int Symp New
Ideas, New Paradigms, Reflections Program Softw; 2019, p. 154--178.
\url{https://doi.org/10.1145/3359591.3359737}.

\bibitem{lewis2020rag}
Lewis P, Perez E, Piktus A, Petroni F, Karpukhin V, Goyal N, et al.
Retrieval-augmented generation for knowledge-intensive NLP tasks. Adv Neural
Inf Process Syst 2020;33:9459--74.
\url{https://arxiv.org/abs/2005.11401}.

\bibitem{yao2023react}
Yao S, Zhao J, Yu D, Du N, Shafran I, Narasimhan K, et al. ReAct: Synergizing
reasoning and acting in language models. In: Int Conf Learn Represent; 2023.
\url{https://arxiv.org/abs/2210.03629}.

\bibitem{yang2024sweagent}
Yang J, Jimenez CE, Wettig A, Lieret K, Yao S, Narasimhan K, et al. SWE-agent:
Agent-computer interfaces enable automated software engineering. Adv Neural Inf
Process Syst 2024;37.
\url{https://arxiv.org/abs/2405.15793}.

\bibitem{agentwiki}
Chu Z. Agent Wiki: A local-first evidence-closed Markdown wiki for LLM agents.
Version 0.2.1 [software]. Zenodo; 2026 Jul 15.
\url{https://doi.org/10.5281/zenodo.21252445}.

\end{thebibliography}

\end{document}
